\section{Módulo de Gastos}
Al iniciar sesión, el usuario aterriza en el módulo de "Gastos". Esta sección cuenta con dos pestañas de navegación, abriendo por defecto la vista de \textbf{Análisis}.

Este módulo opera además como \textbf{libro mayor} de la aplicación: recibe tanto los gastos registrados manualmente como los generados de forma automática por los módulos de Cuotas y Suscripciones.

\subsection{Dashboard de Análisis}
La pestaña de análisis está dedicada a la visualización de métricas financieras e integra actualmente dos componentes gráficos:

\begin{enumerate}
    \item \textbf{Evolución del Gasto (Gráfico de Líneas):}
    Proyecta los gastos en una escala temporal mensual. Incluye una herramienta de control financiero que permite al usuario ajustar un umbral o límite de presupuesto para monitorear posibles sobregiros, dibujado como línea de referencia sobre la serie.

    \item \textbf{Distribución del Gasto (Diagrama de Anillo):}
    Implementado como un \textit{donut} interactivo: al seleccionar un segmento, o su entrada en la leyenda, éste se destaca y el centro del anillo muestra su porcentaje y su valor.

    El componente responde a \textbf{dos preguntas distintas}, alternables mediante un conmutador, porque no son equivalentes:
    \begin{itemize}
        \item \textbf{Monto:} cuánto dinero se gastó con cada categoría.
        \item \textbf{Frecuencia:} cuántos movimientos se registraron con cada una.
    \end{itemize}

El gráfico admite además \textbf{dos ejes de agrupación} distintos, ya que un mismo gasto responde a dos preguntas que no conviene mezclar:
    \begin{itemize}
        \item \textbf{Método:} con qué medio se pagó (efectivo, débito, crédito, transferencia).
        \item \textbf{Tipo:} de dónde proviene el gasto (directo, cuotas o suscripciones).
    \end{itemize}

    Un cobro de un servicio cargado a la tarjeta de crédito es \textit{Crédito} en el primer eje y \textit{Suscripciones} en el segundo; agruparlos en una sola dimensión perdería ambas lecturas. Los movimientos cuyo medio de pago no se ha declarado se agrupan bajo una categoría propia y se informan al pie del gráfico.
\end{enumerate}
